\graphicspath{{../01Introduction/pics/}}

\chapter{Introduction}\label{ch:Introduction}

\begin{quoting}
	This quantum business is so incredibly important and difficult that everyone should busy himself with it. 
	
	
	A. Einstein in a letter to his friend Jakob Laub in 1908, as quoted by A. Wheeler in “The Mystery and The Message Of The Quantum”
\end{quoting}

{\bf Abstract}\hspace{0.2cm} In this chapter.


\lettrine[lines=2]{\color{darkocre}Q}{uantum physics is a century-old branch of physics. Its success} is unparalleled and yet quantum physics is unfinished in one sense: There is no clear and widely adopted consensus on what some of quantum ideas "really mean."

\section{What Is Quantum Physics?}
There are many characterizations of quantum physics.  In essence, quantum physics is the part of physics which focuses on \emph{quantum systems} -- physical systems showing quantum behavior.


\section{Brief Historical Context}
The year 1900 is usually considered the birth year of quantum physics. On December 14 of 1900, at the meeting ????, the German physicist Max Karl Ernst Ludwig Planck presented his theoretical explanation of the spectrum of electromagnetic radiation emitted by hot bodies.  

1802 
\subsection{Three Ages of Quantum}
The evolution of quantum science and technology can be roughly divided into three stages: "Old quantum physics", modern quantum physics, and information-age quantum physics.

\subsubsection*{Old Quantum Physics}

\subsubsection*{Modern Quantum Physics}

\subsubsection*{Information Age}


\section{Who Needs Quantum Physics?}

In October of 1912, Albert Einstein\index{Einstein} wrote in a letter to his physicist
friend Arnold Sommerfeld:
\begin{mybio}{Example of mybio environment}
  I am now exclusively occupied with the problem of gravitation theory
and hope, with the help of a local mathematician friend, to overcome
all the difficulties. One thing is certain, however, that never in my
life have I been quite so tormented. A great respect for mathematics
has been instilled within me, the subtler aspects of which, in my stupidity,
I regarded until now as a pure luxury. Against this problem [of
  gravitation] the original problem of the theory of relativity is
child’s play.
\end{mybio}
In the period from 1905 to 1916 Einstein was feverishly working on the
General Theory of Relativity\index{General relativity} -- the next
best theory of gravity since
Newton. The mathematics of general relativity is based on the calculus
of tensors, created by Italian mathematicians Ricci-Curbastro and
Levi-Civita roughly a decade before Einstein started working on the
problem of gravity.


\section{Why is Quantum Physics Hard?}\label{sec:WhyQuantumHard}
Now what are tensors more rigorously? Can we give a short definition
to this concept?
Let us take a look at several examples and see whether they shed
sufficient light. The definitions given below differ from each other,
but they simply convey \emph{the same idea in different ways}.
section{Challenges}
To illustrate the concepts of functions, operators, their structures and
properties, we will be using schematics\index{Schematic} like the one

\subsection{Mathematics}
To illustrate the concepts of functions, operators, their structures and
properties, we will be using 

\subsection{Language}
To illustrate the concepts of functions, operators, their structures and
properties, we will be using schematics\index{Schematic} like the one

The Encyclopedia of Mathematics
\footnote{\url{https://encyclopediaofmath.org/wiki/Tensor_on_a_vector_space}}
provides the following definition:

%\begin{tcolorbox}[colback=white!85!ocre, title=Tensors Definition 1]
\begin{mydef}{Example of mydef environment}\\
  \small
Tensor\index{Tensor} on a vector space $V$ over a field $k$ is an element $t$ of the
vector space
\begin{equation*}
	T^{p,q} (V) = (\otimes^p V)\otimes (\otimes^q V^*)\,,
\end{equation*}
where $V^*=\textrm{Hom}(V, k)$ is the dual space of $V$.
\label{tcb:tensorDefinition1}
\end{mydef}
%\end{tcolorbox}
To understand this defintion we first need to understand what
\emph{vector space}\index{Vector!space} is, what \emph{field} is, what \emph{dual} means,
and what is going on with superscripts and circles (e.g., in
$\otimes^q$).

\subsection{Concepts}
To illustrate the concepts of functions, operators, their structures and
properties, we will be using schematics\index{Schematic} like the one


\section{Quantum Versus Classical}
Sometimes to illustrate mathematical concepts and \emph{relations between them}\index{Relations},
we will use diagrams. Diagrams are helpful in highlighting
some general features of \emph{mathematical structures}\index{Mathematical!structure}.

\begin{figure}[htbp]
  \centering
  \includegraphics[scale=1.0]{defaultFigureTemplate}
  \caption{Diagrams are used to graphically represent sets of objects
    and relationships between them. Arrows can connect (map) elements
    of one set with another. Such mappings may have names: {\bf mlg}
    returns mileage for a given car, {\bf clr} -- color, and {\bf smk}
  determines whether two cars are of the same make.}
  \label{fig:diagrams}
\end{figure}

A particular property of a car-point
can then be represented using an arrow that connects the car-point to
another point in the relevant set. We say that such an arrow \emph{maps}\index{Map}
points of one set into another set. The Figure\ref{fig:diagrams}(b)
shows three maps: $\textrm{\bf mlg}$ gives the mileage for each car from the
set $\Lambda$, $\textrm{\bf clr}$ gives the color for each car, and
$\textrm{\bf smk}$ compares whether two cars have the same make.

%\begin{tcolorbox}[colback=white!85!ocre, title=Problem]
\begin{exercise}
Extend the diagram from the Figure \ref{fig:diagrams}(b), adding a set
of different car makes (e.g., Ford, Toyota, Fiat, etc.) Come up with a
mapping from this set into the Boolean set $B$.
\label{exe:carMakesSet}
\end{exercise}
%\end{tcolorbox}

\section{Quantum Puzzle}
To illustrate the concepts of functions, operators, their structures and
properties, we will be using schematics\index{Schematic} like the one
shown in the Figure \ref{fig:schematicExample}.
\begin{SCfigure}%[htbp]
  %\centering
  \includegraphics[scale=1.0]{defaultFigureTemplate}
  \caption{Schematics can be used to represent functions, operators,
    their compositions and structure.}
  \label{fig:schematicExample}
\end{SCfigure}

A simple schematic element is represented as a box with inputs
and outputs. A box can have a name (label) which describes what the
function does to its input.  The number of inputs and
outputs can vary depending on the complexity of a function.

\

\vspace{1cm}
\section*{Chapter Highlights}
{\setstretch{1.5}\chhc
  \it
\begin{itemize}
\item Natural evolution of mathematical objects from numbers, through
  vectors, leads to tensors.
\item Each successive tier of mathematical object in the progression
  ``numbers, vectors, tensors''  is more abstract and more powerful.
\item Numbers, vectors, and tensors are all conceptually connected.
\end{itemize}
}
